% Source: https://artofproblemsolving.com/community/c3h2681273_number_theory_and_an_inequality
% Stable identity: aops-topic-2681273-q01
Note $\dfrac{xy}{x+y}=\dfrac1{s}$ where $s=\dfrac1x+\dfrac1y$. The condition is $s<\dfrac12$, and minimising $\dfrac{xy}{x+y}$ means maximising $s$ subject to $s<\dfrac12$. By symmetry let $x\le y$.
\begin{itemize}
\item $x\le2$ is impossible, since then $s>\dfrac1x\ge\dfrac12$.
\item $x=3$: we need $\dfrac1y<\dfrac16$, so $y\ge7$ and $s\le\dfrac13+\dfrac17=\dfrac{10}{21}$.
\item $x=4$: we need $\dfrac1y<\dfrac14$, so $y\ge5$ and $s\le\dfrac14+\dfrac15=\dfrac9{20}<\dfrac{10}{21}$.
\item $x\ge5$: $s\le\dfrac25<\dfrac{10}{21}$.
\end{itemize}
So the largest possible $s$ is $\dfrac{10}{21}$, at $(x,y)=(3,7)$, and the minimum value is $\dfrac{xy}{x+y}=\dfrac{21}{10}$:
\[\boxed{\dfrac{21}{10}}.\]
